\documentclass[11pt]{article}
\usepackage{icmlko}
\begin{document}

\title{같은 값을 두 번 내면 산 게 아니다 --- 규약이 통과시킨 거래를
환율이 막는다}
\author{월드모델 실험실}
\date{노트 460 · 2026년 8월 2일}
\maketitle

\begin{abstract}
남은 손잡이는 \texttt{it}(반복 수) 하나였고 $220$ 이 어디서 온 값인지
장부에 없었다. 깊이 $12 \cdot$ lr $0.10$ 에서 훑었다.

\textbf{판은 $220$ 이 봉우리다}($120$ $-0.0035$ · $400$ $-0.0013$ ·
$700$ $-0.0052$). 그런데 옆 칸이 다른 말을 한다 --- \textbf{집 밖은
\texttt{it} 이 늘수록 단조로 내려간다}(KR: $120$ $+0.5575 \to 220$
$+0.5561 \to 400$ $+0.5478 \to 700$ $+0.5341$).

\textbf{깊이는 집 밖을 사고 \texttt{it} 은 집 밖을 판다.} 노트 $408 \sim
410$ 이 한 가족이라 부른 셋이 집 밖에 대해 \emph{부호가 반대}다.

그래서 $120$ 을 제대로 쟀다(씨앗 열 짝지어).

\begin{center}
\begin{tabular}{lrr}
\hline
계기 & \texttt{it} $220 \to 120$ & 95\% 구간 \\
\hline
판 & $-0.0034$ & $[-0.0041, -0.0027]$ \\
KR 만화 & $\mathbf{+0.0029}$ & $\mathbf{[+0.0016, +0.0042]}$ \\
비게임 앱 & $+0.0008$ & $[-0.0001, +0.0017]$ \\
\hline
\end{tabular}
\end{center}

노트 426 거울의 조건이 \textbf{셋 다 찬다} --- ① 같은 부호 ② KR 구간이
$0$ 을 안 뭄 ③ 판 손해 $0.0034$ 가 $2\times$짝SD($0.0034$) 안.
\textbf{규약대로면 넣어야 한다.}

\textbf{안 넣는다.} 그리고 규약을 고친다.
\end{abstract}

\begin{figure}[t]
\centering
\includegraphics[width=\textwidth]{figs/f1.pdf}
\caption{\textbf{가} --- \texttt{it} 을 늘리면 판은 봉우리를 지나 내려가고
집 밖은 처음부터 계속 내려간다. \textbf{나} --- 지금까지의 거래를 환율로
줄 세웠다. 새 바닥은 $1$ 이다.}
\end{figure}

\section{환율 $0.54$ --- 낸 것보다 적게 산다}

노트 459 가 환율(집 밖 평균 이득 $\div$ 판 손해)을 적기 시작했다.
\texttt{it} $120$ 의 환율은 $0.00185 / 0.0034 = \mathbf{0.54}$ 다.

\textbf{판 rho 와 집 밖 rho 는 같은 단위다.} 둘 다 스피어만이다. 환율이
$1$ 아래라는 것은 \textbf{내놓은 rho 보다 적은 rho 를 사 온다}는 뜻이다.
어느 쪽을 더 치느냐를 정하기 전에도, \emph{같은 저울에 올려 덜 받는
거래}인 것은 분명하다.

\begin{center}
\begin{tabular}{lrrl}
\hline
변경 & 판 손해 & 집 밖 이득 & 환율 \\
\hline
\texttt{OOC\_DROP} (노트 441) & $0.0000$ & $+0.3681$ & $\infty$ \\
깊이 $6 \to 12$ (노트 457) & $0.0032$ & $+0.0063$ & $1.95$ \\
lr $0.10 \to 0.06$ (노트 459) & $0.0025$ & $+0.0016$ & $0.62$ \\
\textbf{\texttt{it} $220 \to 120$} & $0.0034$ & $+0.0019$ & $\mathbf{0.54}$ \\
\hline
\end{tabular}
\end{center}

\section*{규약 수정 --- 환율 바닥 $1$}

노트 458 이 \emph{거래 한도가 없다}고 적었고, 노트 459 는 \emph{조건 ②가
대부분을 막으니 한도 대신 기록만 하자}고 걸었다. \textbf{그 내기가 한
노트 만에 졌다.} \texttt{it} $120$ 은 조건 ②를 통과한다 --- KR 구간이
깨끗하게 $0$ 을 안 문다. 값이 작아서가 아니라 \emph{정밀해서} 통과했다.

\begin{quote}
\textbf{조항} --- 판을 잃고 집 밖을 사는 변경은 \textbf{환율이 $1$ 이상}
이라야 넣는다. 판 rho 와 집 밖 rho 는 같은 단위이므로, 환율 $1$ 아래는
\emph{같은 저울에서 덜 받는 거래}다. 노트 426 거울의 조건 ①②③에
\textbf{④ 환율 $\ge 1$} 을 더한다.
\end{quote}

이 조항은 지금까지의 판정을 하나도 안 뒤집는다 --- \texttt{OOC\_DROP}
($\infty$)과 깊이 $12$($1.95$)는 그대로 서고, lr $0.06$($0.62$)은 이미
조건 ②로 기각됐으며, \texttt{it} $120$($0.54$)만 새로 막힌다.
\textbf{바꾸는 규약이 과거를 다시 쓰지 않는지 확인하는 것}이 이 문단의
용도다.

\section{덤 --- 노트 428 이 두 번째로 어긋난다}

노트 428 은 \emph{관문(되뽑기)이 용량 변경을 아래로 어림한다}고 했다.
노트 457 에서 이미 방향이 어긋났고(관문 $-0.0028$ 대 전 자료 $-0.0032$),
여기서 더 크게 어긋난다 --- \textbf{관문 $-0.0004$($5/12$) 대 전 자료
$-0.0034$}. 관문이 여덟 배 너그럽다.

용량을 \emph{줄이는} 변경에서도 같은 방향이니 ``되뽑기가 학습을 $63\%$ 로
줄여 깊은 나무에 불리하다''는 설명으로는 안 맞는다. \textbf{노트 428 의
어림은 이 챔피언에서 더 쓰지 않는다} --- 판정은 전 자료 짝씨앗 수치로
한다. (왜 어긋나는지는 안 쟀다.)

\section*{남은 것}

\begin{itemize}
\item \textbf{왜 \texttt{it} 은 집 밖을 파나.} 깊이는 상호작용을 늘리고
\texttt{it} 은 잔차를 더 깎는다. \emph{잔차를 깎는 것이 수집 고유의
무늬를 외우는 일}이라면 그게 답이다 --- 학습 안 적합도를 \texttt{it}
별로 재면 갈린다.
\item \textbf{용량 셋이 다 제자리다.} 깊이 $12$ · lr $0.10$ · \texttt{it}
$220$ · K $32$. 이 챔피언의 용량 쪽은 닫혔다고 적는다.
\item \textbf{환율 바닥을 정할 근거가 더 필요하다.} $1$ 은 ``같은 단위''
라는 것 말고 근거가 없다. 판과 집 밖의 값어치를 실제로 재려면
\emph{배포가 어느 쪽을 쓰는지}를 봐야 한다.
\end{itemize}

\end{document}
